\Needspace{12\baselineskip}
\begin{center}
\textit{Division of classes into orders.}
\end{center}

\Needspace{8\baselineskip}
\Article{226}
We call a form $(a,b,c)$ \textit{primitive} if the numbers $a,b,c$ share no nontrivial divisor; otherwise it is called \textit{derived}. More precisely, if the greatest common divisor of the numbers $a,b,c$ is set $=m$, the form $(a,b,c)$ is \textit{derived from the primitive form} $\left(\frac{a}{m},\frac{b}{m},\frac{c}{m}\right)$. This definition immediately shows that all forms whose determinant is divisible by no square (except $1$) must be primitive. Further, Article 161 shows that if a given class of forms with determinant $D$ contains a primitive form, all forms of this class are primitive; in that case
\SourceLeaf{232}
the class itself is called \textit{primitive}. It is also clear that if a form $F$ with determinant $D$ is derived from a primitive form $f$ with determinant $\frac{D}{mm}$, and the classes containing the forms $F$, $f$ are $K$, $k$, respectively, then all forms in the class $K$ are derived from the primitive class $k$. In this case we therefore call the class $K$ itself \textit{derived from the primitive class $k$}.

If $(a,b,c)$ is a primitive form, but $a$, $c$ are not both even (\textit{i.e.}, both are odd, or at least one is), it is easy to see that not only $a$, $b$, $c$, but also $a$, $2b$, $c$ share no nontrivial divisor. In this case the form $(a,b,c)$ is called \textit{properly primitive}, or simply a \textit{proper form}. But if $(a,b,c)$ is a primitive form and the numbers $a$, $c$ are both even, the numbers $a$, $2b$, $c$ clearly have the common divisor $2$ (which is also their greatest common divisor), and $(a,b,c)$ is called an \textit{improperly primitive form}, or simply an \textit{improper form}\HistoricalNote{*)}{We chose the terms \textit{properly} and \textit{improperly} here because no more suitable ones came to mind. We mention this so that no one will look for a hidden connection between this meaning and the one we have used since Article 157: there is none. In any case, there is certainly no need to fear ambiguity here.}. In this case $b$ must be odd (otherwise $(a,b,c)$ would not be primitive); so $bb\equiv1\ (\mathrm{mod.}\ 4)$, and since $ac$ is divisible by $4$, the determinant $bb-ac\equiv1\ (\mathrm{mod.}\ 4)$. Improper forms therefore occur only for a determinant of the form $4n+1$ if positive, or of the form $-(4n+3)$ if negative. --- Article 161 also makes it clear that if a given class contains a properly primitive form, all forms of this class are properly primitive; conversely, a class containing an improperly primitive form consists only of improperly primitive forms. Thus in the first case the class itself is called \textit{properly primitive}, or simply \textit{proper}; in the second, \textit{improperly primitive}, or \textit{improper}. For \textit{example}, among the positive classes of forms with determinant $-235$, six are proper, namely those represented by $(1,0,235)$, $(4,1,59)$, $(4,-1,59)$, $(5,0,47)$, $(13,5,20)$, $(13,-5,20)$, and there are just as many among the negative classes; two in each group are improper. --- The classes of forms with determinant $79$ (a number of the form $4n+3$) are all proper.

If a form $(a,b,c)$ is derived, specifically from the primitive form $\left(\frac{a}{m},\frac{b}{m},\frac{c}{m}\right)$, the latter can be either properly or improperly primitive. In the first case $m$ is also the greatest common divisor of the numbers $a$, $2b$, $c$; in the second their greatest common divisor is $2m$. This explains the distinction between a \textit{form derived from a properly primitive form} and a \textit{form derived from an improperly primitive form}; and also (since by Article 161 all forms in the same class behave alike in this respect)
\SourceLeaf{233}
between a \textit{class derived from a properly primitive class} and a \textit{class derived from an improperly primitive class}.

These distinctions give us a first foundation for dividing all classes of forms with a given determinant into different \textit{orders}. We place two classes represented by the forms $(a,b,c)$, $(a',b',c')$ in the \textit{same order} if \textit{both} the numbers $a$, $b$, $c$ have the same greatest common divisor as $a'$, $b'$, $c'$, \textit{and} $a$, $2b$, $c$ have the same greatest common divisor as $a'$, $2b'$, $c'$. If either or both of these conditions fail, the classes belong to \textit{different orders}. It follows at once that all properly primitive classes form one order, and all improperly primitive classes another. If $mm$ is a square dividing the determinant $D$, the classes derived from properly primitive classes with determinant $\frac{D}{mm}$ form a separate order, and those derived from improperly primitive classes with determinant $\frac{D}{mm}$ another, etc. If $D$ is divisible by no square (except $1$), there are no orders of derived classes. Thus there is either just one order (when $D\equiv2$ or $3$ modulo $4$), namely the order of properly primitive classes, or two (when $D\equiv1\ (\mathrm{mod.}\ 4)$), namely the order of properly primitive classes and the order of improperly primitive classes. The principles of combinatorics readily give the following general rule. Suppose $D=D'2^{2\mu}a^{2\alpha}b^{2\beta}c^{2\gamma}\ldots.$, where $D'$ has no square factor and $a$, $b$, $c$ etc. are distinct odd primes (any number can be put in this form by setting $\mu=0$ when $D$ is not divisible by $4$, and setting $\alpha$, $\beta$, $\gamma$ etc. all $=0$, or, equivalently, omitting the factors $a^{2\alpha}$, $b^{2\beta}$, $c^{2\gamma}$ etc., when $D$ is divisible by no odd square). Then there are either
\[
(\mu+1)(\alpha+1)(\beta+1)(\gamma+1)\ldots
\]
orders, namely when $D'\equiv2$ or $3\ (\mathrm{mod.}\ 4)$, or
\[
(\mu+2)(\alpha+1)(\beta+1)(\gamma+1)\ldots
\]
orders when $D'\equiv1\ (\mathrm{mod.}\ 4)$. But we omit the proof of this rule, since it is neither difficult nor particularly necessary here.

\textit{Example} 1. For $D=45=5.3^2$ there are six classes, represented by $(1,0,-45)$, $(-1.\,0,45)$, $(2,1,-22)$, $(-2,1,22)$, $(3.\,0,-15)$, $(6,3,-6)$. They are divided into four orders: Order I contains the two proper classes represented by $(1,0,-45)$, $(-1,0,45)$; Order II contains the two improper classes\SourceLeaf{234}, represented by $(2,1,-22)$, $(-2,1,22)$; Order III contains one class derived from a proper class with determinant $5$, namely the one represented by $(3,0,-15)$; Order IV consists of one class derived from an improper class with determinant $5$, namely the one represented by $(6,3,-6)$.

\textit{Example} 2. The positive classes with determinant $-99=-11.3^2$ are divided into four orders. Order I comprises the following properly primitive classes\HistoricalNote{*)}{For brevity, we use representative forms in place of the classes they represent.}: $(1,0,99)$, $(4,1,25)$. $(4,-1,25)$, $(5,1,20)$, $(5,-1,20)$, $(9,0,11)$; Order II contains the improper classes $(2,1,50)$, $(10,1,10)$; Order III, the classes derived from proper classes with determinant $-11$, $(3,0,33)$, $(9,3,12)$, $(9,-3,12)$; Order IV, the single class derived from an improper class with determinant $-11$, $(6,3,18)$. --- The negative classes with this determinant can be divided into orders in exactly the same way.

We note that \textit{opposite classes always belong to the same order}; the reason for this theorem is easy to see.

\Needspace{8\baselineskip}
\Article{227}
Among these different orders, the order of properly primitive classes especially deserves the closest attention. Each derived class comes from a certain primitive class (with a smaller determinant), and considering the latter usually gives what concerns the former directly. Below we will show that every improperly primitive class is, in a similar way, associated, so to speak, with either a single properly primitive class or three (with the same determinant). Further, for negative determinants we may pass over the negative classes, since a certain positive class always corresponds to each of them. To understand the nature of properly primitive classes more deeply, we will first explain a certain essential distinction by which the whole order of proper classes can be subdivided into several \textit{genera}. Since we have not yet touched on this very important subject, we must take it up from the beginning.

\Needspace{12\baselineskip}
\begin{center}
\textit{Division of orders into genera.}
\end{center}

\Needspace{8\baselineskip}
\Article{228}
\textsc{Theorem.} \textit{Any properly primitive form $F$ represents infinitely many numbers not divisible by any given prime $p$.}

\textit{Proof.} If the form $F=axx+2bxy+cyy$, clearly $p$ cannot divide all three numbers $a$, $2b$, $c$ at once. Now if $a$ is not divisible by $p$\SourceLeaf{235}, it is clear that taking for $x$ any number not divisible by $p$, and for $y$ a number divisible by $p$, makes the value of the form $F$ not divisible by $p$. If $c$ is not divisible by $p$, the same result follows by giving $x$ a divisible value and $y$ a nondivisible one. Finally, if both $a$ and $c$ are divisible by $p$, so that $2b$ is not, the form $F$ takes a value not divisible by $p$ when both $x$ and $y$ are given any values not divisible by $p$. \textit{Q.E.D.}

The theorem clearly also holds for \textit{improperly primitive} forms, provided $p=2$ does not hold.

Since several conditions of this kind can hold at once, with the same number divisible by certain given primes and not divisible by others (see Article 32), it is easy to see that the numbers $x$, $y$ can be chosen in infinitely many ways so that the primitive form $axx+2bxy+cyy$ takes a value not divisible by any number of given primes. Only $2$ must be excluded from those primes when the form is improperly primitive. Thus we can state the theorem more generally as follows: \textit{Any primitive form represents infinitely many numbers relatively prime to any given number (which must be odd when the form is improperly primitive).}

\Needspace{8\baselineskip}
\Article{229}
\textsc{Theorem.} \textit{Let $F$ be a primitive form with determinant $D$, and $p$ a prime dividing $D$. Then the numbers not divisible by $p$ that can be represented by the form $F$ agree in this respect: either they are all quadratic residues of $p$, or they are all nonresidues.}

\textit{Proof.} Let $F=(a,b,c)$; let $m$, $m'$ be any two numbers not divisible by $p$ that can be represented by the form $F$, namely
\[
m=agg+2bgh+chh,\qquad m'=ag'g'+2bg'h'+ch'h'
\]
Then
\[
mm'=(agg'+b(gh'+hg')+chh')^2-D(gh'-hg')^2
\]
so $mm'$ is congruent to a square modulo $D$, and therefore also modulo $p$, \textit{i.e.}, $mm'$ is a quadratic residue of $p$. It follows that either both $m$, $m'$ are quadratic residues of $p$, or both are nonresidues. \textit{Q.E.D.}

\SourceLeaf{236}
Similarly, when the determinant $D$ is divisible by $4$, we prove that all odd numbers representable by $F$ are either all $\equiv1$. or all $\equiv3\ (\mathrm{mod.}\ 4)$. For in this case the product of two such numbers is always a quadratic residue of $4$, and so $\equiv1\ (\mathrm{mod.}\ 4)$; hence either both are $\equiv1$, or both $\equiv3$.

Finally, when $D$ is divisible by $8$, the product of any two odd numbers representable by $F$ is a quadratic residue of $8$, and therefore $\equiv1\ (\mathrm{mod.}\ 8)$. Thus in this case the odd numbers representable by $F$ are either all $\equiv1$, or all $\equiv3$, or all $\equiv5$, or all $\equiv7\ (\mathrm{mod.}\ 8)$.

Thus, for \textit{example}, since the form $(10,3,17)$ represents the number $10$, which is a nonresidue of $7$, all numbers not divisible by $7$ that this form represents are nonresidues of $7$. --- Since $-3$ is representable by the form $(-3,1,49)$ and modulo $4$ is $\equiv1$, all odd numbers representable by this form behave in the same way.

If it were needed for our present purpose, we could easily prove that the numbers representable by the form $F$ have no such fixed relation to any prime not dividing $D$: both residues and nonresidues of any prime not dividing $D$ can be represented by the form $F$ without distinction. For the numbers $4$ and $8$, on the other hand, something similar also holds in other cases, which we cannot pass over.

I. \textit{When the determinant $D$ of the primitive form $F$ is $\equiv3\ (\mathit{mod.}\ 4)$, the odd numbers representable by the form $F$ are either all $\equiv1$, or all $\equiv3\ (\mathit{mod.}\ 4)$.} For if $m$, $m'$ are two numbers representable by $F$, their product $mm'$ can be put in the form $pp-Dqq$ as above. If both $m$, $m'$ are odd, one of the numbers $p$, $q$ must be even and the other odd, so one of the squares $pp$, $qq$ is $\equiv0$ and the other $\equiv1\ (\mathrm{mod.}\ 4)$. It follows easily that $pp-Dqq$ is certainly $\equiv1\ (\mathrm{mod.}\ 4)$, and so either both $m$, $m'$ are $\equiv1$, or both $\equiv3\ (\mathrm{mod.}\ 4)$. For \textit{example}, the form $(10,3,17)$ cannot represent odd numbers other than those of the form $4n+1$.

II. \textit{When the determinant $D$ of the primitive form $F$ is $\equiv2\ (\mathit{mod.}\ 8)$, the odd numbers representable by the form $F$ are either partly $\equiv1$ and partly $\equiv7$, or partly $\equiv3$ and partly $\equiv5\ (\mathit{mod.}\ 8)$.} For let $m$, $m'$ be two odd numbers
\SourceLeaf{237}
representable by $F$, so that their product $mm'$ can be put in the form $pp-Dqq$. Since both $m$, $m'$ are odd, $p$ must be odd (because $D$ is even), and so $pp\equiv1\ (\mathrm{mod.}\ 8)$; but $qq$ is either $\equiv0$, $\equiv1$, or $\equiv4$, and therefore $Dqq$ is either $\equiv0$ or $\equiv2$. Hence $mm'=pp-Dqq$ is either $\equiv1$ or $\equiv7\ (\mathrm{mod.}\ 8)$. Thus if $m$ is either $\equiv1$ or $\equiv7$, $m'$ is also either $\equiv1$ or $\equiv7$; but if $m$ is either $\equiv3$ or $\equiv5$, $m'$ is also either $\equiv3$ or $\equiv5$. For \textit{example}, all odd numbers representable by the form $(3,1,5)$ are either $\equiv3$ or $\equiv5\ (\mathrm{mod.}\ 8)$, and no numbers of the form $8n+1$ or $8n+7$ can be represented by that form.

III. \textit{When the determinant $D$ of the primitive form $F$ is $\equiv6\ (\mathit{mod.}\ 8)$, this form can represent odd numbers only in the two cases $\equiv1$ and $\equiv3\ (\mathit{mod.}\ 8)$, or only in the two cases $\equiv5$ and $\equiv7\ (\mathit{mod.}\ 8)$.} Anyone can easily work out a proof entirely similar to the preceding one (in II). --- Thus, for \textit{example}, the only odd numbers representable by the form $(5,1,7)$ are those either $\equiv5$ or $\equiv7\ (\mathrm{mod.}\ 8)$.

\Needspace{8\baselineskip}
\Article{230}
Thus all numbers representable by a given primitive form $F$ with determinant $D$ have a fixed relation to each prime divisor of $D$ (by which they themselves are not divisible). The odd numbers representable by $F$ also have, in certain cases, a fixed relation to the numbers $4$ and $8$: namely to $4$ whenever $D$ is either $\equiv0$ or $\equiv3\ (\mathrm{mod.}\ 4)$, and to $8$ whenever $D$ is either $\equiv0$, $\equiv2$, or $\equiv6\ (\mathrm{mod.}\ 8)$\HistoricalNote{*)}{For determinants divisible by $8$, the relation to the number $4$ can be left out, since in this case it is already contained in the relation to $8$.}. We call such a relation to each of these numbers a \textit{character}, or an \textit{individual character}, of the form $F$, and express it as follows. When only quadratic residues of the prime $p$ can be represented by the form $F$, we assign it the character $Rp$, and in the opposite case the character $Np$. Similarly, we write $1,4$ when the form $F$ cannot represent any odd numbers except those that are $\equiv1\ (\mathrm{mod.}\ 4)$. This makes it immediately clear which characters the signs $3,4$; $1,8$; $3,8$; $5,8$; $7,8$ express. Finally, to forms that represent only odd numbers that modulo
\SourceLeaf{238}
$8$ are either $\equiv1$ or $\equiv7$, we assign the character $1\ \mathit{and}\ 7,8$. The meaning of the characters $3\ \mathit{and}\ 5,8$; $1\ \mathit{and}\ 3,8$; $5\ \mathit{and}\ 7,8$ then follows by itself.

The individual characters of a given primitive form $(a,b,c)$ with determinant $D$ can always be found from at least one of the numbers $a$, $c$ (both of which are clearly representable by that form). For whenever $p$ is a prime divisor of $D$, at least one of the numbers $a$, $c$ is certainly not divisible by $p$. If both were divisible by $p$, then $p$ would also divide $bb(=D+ac)$. and therefore $b$ as well; \textit{i.e.}, the form $(a,b,c)$ would not be primitive. Similarly, whenever the form $(a,b,c)$ has a fixed relation to the number $4$ or $8$, at least one of the numbers $a$, $c$ is certainly odd, and this determines that relation. For \textit{example}, the character of the form $(7,0,23)$ with respect to the number $23$ follows from the number $7$ and is $N23$; the character of the same form with respect to $7$ is obtained from the number $23$, namely $R7$; finally, this form's character with respect to $4$, namely $3,4$, can be found from either $7$ or $23$.

Since all numbers representable by a form $F$ in a class $K$ are also representable by every other form of that class, each character of the form $F$ clearly also belongs to all the other forms in this class. We may therefore regard them as characters of the whole class. Thus the individual characters of any given primitive class are found from its representative form. Opposite classes always have all their characters in common.
